\section{Recitation 2: Boolean Expressions}
\label{sec:rec2}

This recitation builds the course's first interpreter and first program transformation, for a tiny language of Boolean expressions. Everything is in plain Racket, with structs and \code{match}.

\subsection{The Language and Its Representation}

The grammar of \emph{bexp}:
\[
\langle\textit{bexp}\rangle \;::=\; \langle\textit{boolean}\rangle \;\mid\; \langle\textit{bexp}\rangle \land \langle\textit{bexp}\rangle \;\mid\; \langle\textit{bexp}\rangle \lor \langle\textit{bexp}\rangle \;\mid\; \neg \langle\textit{bexp}\rangle
\]
One struct per alternative of the grammar represents it in Racket:
\begin{lstlisting}
(struct leaf-node (arg))
(struct or-node (left right))
(struct and-node (left right))
(struct not-node (arg))

(define bexp-bool-1 (and-node (leaf-node #t) (leaf-node #f)))
(define bexp-sym-1 (and-node (leaf-node 'x) (leaf-node 'y)))
\end{lstlisting}
A value built from these structs is an \emph{abstract syntax tree}: the tree structure replaces parentheses and precedence. Note that the structs are \emph{not} transparent and their fields are unchecked: \code{bexp-sym-1} has symbols, not booleans, at its leaves, and Racket accepts it.

\begin{exercise}[Syntactic equality]
Write \code{syn-equal?}, which checks whether two bexps are the same tree.
\end{exercise}
\begin{solution}
Because the structs are opaque, \code{equal?} would compare them by identity (\Cref{sec:structs}), so we write structural equality ourselves. \code{match*} matches several values at once, one pattern per value:
\begin{lstlisting}
(define (syn-equal? b1 b2)
  (match* (b1 b2)
    [((and-node l1 r1) (and-node l2 r2)) (and (syn-equal? l1 l2) (syn-equal? r1 r2))]
    [((or-node l1 r1) (or-node l2 r2))   (and (syn-equal? l1 l2) (syn-equal? r1 r2))]
    [((not-node a1) (not-node a2))       (syn-equal? a1 a2)]
    [((leaf-node v1) (leaf-node v2))     (equal? v1 v2)]
    [(_ _) #f]))
\end{lstlisting}
One clause per variant where both sides agree, and a final catch-all for mismatched shapes.
\end{solution}

\begin{pitfall}[A missing variant hides behind a catch-all]
If the \code{not-node} clause were left out, two \code{not-node}s would fall through to \rkt|[(_ _) #f]|, and \code{syn-equal?} would say a negation is not equal to itself. Nothing would complain, and every later test comparing expressions with a negation would fail even with a correct optimizer. With a catch-all clause present, check the list of variants, not just the clauses you were given.
\end{pitfall}

\subsection{An Interpreter for bexp}

\begin{exercise}[Interpreter]
Write \code{evaluate-bexp}, which computes the Boolean value an expression represents.
\end{exercise}
\begin{solution}
One clause per variant, and a recursive call wherever a variant contains a sub-expression:
\begin{lstlisting}
(define (evaluate-bexp b)
  (match b
    [(leaf-node v)    v]
    [(or-node b1 b2)  (or (evaluate-bexp b1) (evaluate-bexp b2))]
    [(and-node b1 b2) (and (evaluate-bexp b1) (evaluate-bexp b2))]
    [(not-node b1)    (not (evaluate-bexp b1))]))

(check-expect (evaluate-bexp (leaf-node #t)) #t)
(check-expect (evaluate-bexp (and-node (leaf-node #t) (leaf-node #f))) #f)
(check-expect (evaluate-bexp (and-node (leaf-node #f) (leaf-node #t))) #f)
\end{lstlisting}
The third test swaps the arguments of the second; testing both orders is cheap and catches an implementation that looks at only one side.
\end{solution}

This is the pattern of every interpreter in the course: our language's $\land$ is implemented by Racket's \code{and}, its $\lor$ by \code{or}, its $\neg$ by \code{not}. The \emph{meaning} of each construct is given by the clause that handles it.

\subsubsection{Missing cases} What happens with an incomplete \code{match}?
\begin{lstlisting}
(define (evaluate-bexp-m b)
  (match b
    [(leaf-node v) v]))
\end{lstlisting}
Nothing complains when it is defined. The error appears only when it is called on a value that no clause matches: \rkt|(evaluate-bexp-m bexp-bool-1)| raises \code{match: no matching clause}. Racket checks patterns at run time, so a missing case is found only if some test happens to reach it. \Cref{sec:rec3} redoes this in \code{plai-typed}, where the same mistake is rejected before the program runs.

\subsection{An Optimization Pass: De Morgan}

De Morgan's laws and double-negation elimination state the equivalences
\[
\neg(P \land Q) = \neg P \lor \neg Q, \qquad \neg(P \lor Q) = \neg P \land \neg Q, \qquad \neg\neg P = P.
\]
Using them left to right, an optimizer can eliminate redundant negations and push every remaining negation down to the leaves. The examples use symbolic leaves:
\begin{center}
\begin{tabular}{@{}lll@{}}
\toprule
name & before & after \\
\midrule
\code{bexp-1} & $\neg(\neg P \lor Q)$ & $P \land \neg Q$ \\
\code{bexp-2} & $\neg\neg P \lor \neg Q$ & $P \lor \neg Q$ \\
\code{bexp-3} & $P \lor \neg Q$ & $P \lor \neg Q$ (already optimal) \\
\bottomrule
\end{tabular}
\end{center}
\begin{lstlisting}
(define bexp-1 (not-node (or-node (not-node (leaf-node 'P)) (leaf-node 'Q))))
(define bexp-2 (or-node (not-node (not-node (leaf-node 'P))) (not-node (leaf-node 'Q))))
(define bexp-3 (or-node (leaf-node 'P) (not-node (leaf-node 'Q))))

(define bexp-1-opt (and-node (leaf-node 'P) (not-node (leaf-node 'Q))))
(define bexp-2-opt (or-node (leaf-node 'P) (not-node (leaf-node 'Q))))
(define bexp-3-opt (or-node (leaf-node 'P) (not-node (leaf-node 'Q))))
\end{lstlisting}

\subsubsection{What does it mean for an optimization to be correct?} Before writing the optimizer, decide how to check it. An optimization is correct when it does two things:
\begin{enumerate}
  \item it \emph{preserves semantics}: the output evaluates to the same value as the input;
  \item it actually \emph{optimizes}: the output has the promised shape, here every negation sitting directly on a leaf.
\end{enumerate}
Both become executable checks:
\begin{lstlisting}
(define (preserves-sem? unopt)
  (equal? (evaluate-bexp unopt) (evaluate-bexp (demorgan-opt unopt))))

(define (is-optimizing? opt)
  (define (only-not-at-leaves? bexp)
    (match bexp
      [(leaf-node v) #t]
      [(or-node b1 b2) (and (only-not-at-leaves? b1) (only-not-at-leaves? b2))]
      [(and-node b1 b2) (and (only-not-at-leaves? b1) (only-not-at-leaves? b2))]
      [(not-node (leaf-node v)) #t]
      [(not-node _) #f]))
  (only-not-at-leaves? opt))

(define (is-correct? unopt)
  (and (preserves-sem? unopt) (is-optimizing? (demorgan-opt unopt))))
\end{lstlisting}
\code{is-optimizing?} relies on \code{match} trying clauses in order: \rkt|(not-node (leaf-node v))| is tried first, so the catch-all \rkt|(not-node _)| sees only negations of compound expressions, which mean the pass left work undone. Note also that \code{preserves-sem?} takes only the input and calls \code{demorgan-opt} itself, so it tests the optimizer rather than a hand-written answer.

\subsubsection{The optimizer} What to do with a negation depends entirely on what is \emph{inside} it, so a single \rkt|(not-node b1)| clause cannot work. Nested patterns split it into the four shapes the argument can have, one per variant. This is what the handout means by ``you will need more cases than in \code{evaluate-bexp}'':
\begin{lstlisting}
(define (demorgan-opt b)
  (match b
    ; Standard Stuff
    [(leaf-node v)    (leaf-node v)]
    [(or-node b1 b2)  (or-node (demorgan-opt b1) (demorgan-opt b2))]
    [(and-node b1 b2) (and-node (demorgan-opt b1) (demorgan-opt b2))]
    [(not-node (leaf-node l)) (not-node (leaf-node l))]
    ; De-Morgan Stuff
    [(not-node (and-node b1 b2))
     (or-node (demorgan-opt (not-node b1)) (demorgan-opt (not-node b2)))]
    [(not-node (or-node b1 b2))
     (and-node (demorgan-opt (not-node b1)) (demorgan-opt (not-node b2)))]
    [(not-node (not-node b1)) (demorgan-opt b1)]))
\end{lstlisting}
\begin{itemize}
  \item \textbf{A negated leaf} is already as far down as a negation can go; keep it.
  \item \textbf{De Morgan.} After rewriting $\neg(b_1 \land b_2)$ to $\neg b_1 \lor \neg b_2$, the new negations may themselves sit on top of an \code{and}, an \code{or}, or a \code{not}, so they are optimized again: the recursive calls are on \rkt|(not-node b1)|, not on \code{b1}. Calling \rkt|(demorgan-opt b1)| would drop the negation and change the meaning.
  \item \textbf{Double negation.} The two negations cancel, and the result is \rkt|(demorgan-opt b1)|, not just \code{b1}, because \code{b1} may still contain negations to push down.
\end{itemize}

\begin{figure}[htbp]
\centering
\begin{tikzpicture}[level distance=9mm, sibling distance=14mm, every node/.style={ast}]
  \node[asth] {not}
    child {node {or}
      child {node[asth] {not} child {node[astl] {P}}}
      child {node[astl] {Q}}};
  \node[draw=none, fill=none, font=\small] at (0,-3.6) {$\neg(\neg P \lor Q)$};
\end{tikzpicture}
\hspace{0.8cm}
\begin{tikzpicture}
  \draw[arr] (0,0) -- (1.6,0) node[midway, above, lbl] {\code{demorgan-opt}};
  \node at (0,-1.6) {};
\end{tikzpicture}
\hspace{0.8cm}
\begin{tikzpicture}[level distance=9mm, sibling distance=14mm, every node/.style={ast}]
  \node {and}
    child {node[astl] {P}}
    child {node[asth] {not} child {node[astl] {Q}}};
  \node[draw=none, fill=none, font=\small] at (0,-3.6) {$P \land \neg Q$};
\end{tikzpicture}
\caption{\code{bexp-1} before and after the pass. Negations (yellow) end up only directly above leaves.}
\label{fig:demorgan}
\end{figure}

For \code{bexp-1} (\Cref{fig:demorgan}):
\begin{lstlisting}[style=repl]
(demorgan-opt (not (or (not P) Q)))
---> (and-node (demorgan-opt (not (not P))) (demorgan-opt (not Q)))   ; not over or
---> (and-node (demorgan-opt P) (not Q))                              ; not-not cancels
---> (and-node P (not Q))
\end{lstlisting}

\subsubsection{Why it terminates} Every recursive call is on a smaller tree. In the De Morgan cases the argument \rkt|(not-node b1)| is new, but it is smaller than the input \rkt|(not-node (and-node b1 b2))|, which has one more node and all of \code{b2} besides.

\subsection{Testing the Optimizer Two Ways}

The syntactic tests compare against the exact trees written by hand:
\begin{lstlisting}
(check-expect (syn-equal? (demorgan-opt bexp-1) bexp-1-opt) #t)
(check-expect (syn-equal? (demorgan-opt bexp-2) bexp-2-opt) #t)
(check-expect (syn-equal? (demorgan-opt bexp-3) bexp-3-opt) #t)
\end{lstlisting}
The semantic tests check the two correctness properties without writing the expected tree at all:
\begin{lstlisting}
(define bexp-4 (not-node (or-node (not-node (leaf-node #t)) (leaf-node #f))))
(define bexp-5 (or-node (not-node (not-node (leaf-node #t))) (not-node (leaf-node #f))))
(define bexp-6 (or-node (leaf-node #t) (not-node (leaf-node #t))))
(check-expect (is-correct? bexp-4) #t)
(check-expect (is-correct? bexp-5) #t)
(check-expect (is-correct? bexp-6) #t)
\end{lstlisting}

\begin{remark}[Why booleans at the leaves]
\code{bexp-4} to \code{bexp-6} are \code{bexp-1} to \code{bexp-3} with \code{\#t} and \code{\#f} in place of \code{'P} and \code{'Q}. With symbols at the leaves, \code{evaluate-bexp} is meaningless: \rkt|(not 'P)| is \code{\#f} and \rkt|(or 'P #f)| is \code{'P}, because Racket treats every value except \code{\#f} as true. So a semantic check only means something on Boolean leaves.
\end{remark}

\begin{remark}[The handout's table]
An earlier version of the handout listed $\neg(\neg P \land Q) \mapsto P \land Q$ as the first example. That contradicts both the code (\code{bexp-1} is built with \code{or-node}) and De Morgan's law; the posted solution corrects the row to $\neg(\neg P \lor Q) \mapsto P \land \neg Q$. Trust the definitions over the prose, and check the prose by working one example by hand.
\end{remark}
